\abstract{\pic is a young and nearby planetary system hosting an edge-on debris disk and at least two gas-gian planets.
%
Their circumplanetary environments could host exomoons and rings, a detection of which would be highly informative for moon and planet formation theories.
%
A photometric fluctuation of about 4\% was seen towards Beta Pictoris in 1981, indicating the transit of dust in the system, which could be associated with the Hill spheres of the two known planets.}
{We search for the origin of the 1981 event by searching for an analagous event in multi-epoch photometry from 03-02-2017 to 17-01-2023, and look for signs of circumplanetary material in transits of the Hill sphere of \pic c.}
{Observations from the BRITE satellite and the bRING and ASTEP observatories were fitted to a model of the 1981 event to search for a similar event.
%
We searched for a signal consistent with a circumplanetary disk transit using a simple flat disk model.}
{No compelling evidence for the 1981 event was found in photometry between 28-11-2016 (MJD 57720) and 06-01-2023 (MJD 59950).
%
The circumplanetary disk model was limited to transits with an uncertainty in the midpoint of less than 10 days.
%
Due to uncertainty in the orbital solution for c, a search for a disk during the primary transit of MJD 52810 could not be determined, but BRING photometry for the second primary transit at MJD 59413 places an upper limit on the dust content of the Hill sphere of $5\times10^{21}$ grams of material.}
{With no compelling detection of the 1981 event either during the transits of c or other times during the photometric time series, we rule out the hypothesis that this event was related to the Hill spheres of c and b.
%
%Propagated uncertainties in the orbital bundles for planet b prevented detection limits for the Hill sphere of c being determined.
%
Future observations of the next Hill sphere transit in 2028 can be realised with both ground based observatories and with PLATO, whose first long duration science pointing will include Beta Pictoris.
} 
\keywords{planets and satellites: rings – planets and satellites: formation – stars: individual: $\beta \ Pictoris$ }

